%======================================================================
\section{Proofs for Section~\ref{sec:houghton}}\label{app:houghton}
%======================================================================

We use the rules of~\cite[Appendix~C.1]{DouglasMghirbiC}, numbered as there, for unitaries $u,v$ and $\nrm W\le1$: (W2) $\chi(uv)\le\chi(u)+\chi(v)$ and $\chi(u^*)=\chi(u)$; (W4)
$|E(uWu^*)-E(W)|\le2\chi(u)$. Work in the Hilbert--Schmidt space of operators on the Hilbert space $\mathcal H$, where left multiplication acts on the system and right
multiplication on the reference; $\rho$ has finite rank, so $\xi=\sqrt\rho$ is Hilbert--Schmidt and purifies $\rho$, the reference has entropy
$S(\rho)$, and (W2), (W4) and every step below hold verbatim.

\begin{proof}[Proof of Lemma~\ref{lem:sector}]
Write $c_i=\sum_{k\in\Z/3}\omega_3^kP_i^k$ with $\omega_3=e^{2\pi i/3}$, and $Q_{\mathbf a}=P_{i-1}^{a_{i-1}}\cdots P_0^{a_0}$. The $c_k$ with $k<i$ commute, so $\sum_{\mathbf a}Q_{\mathbf a}^*Q_{\mathbf a}=I$,
and $c_i,x_i$ commute with $Q_{\mathbf a}$. Measure $c_0,\dots,c_{N-1}$ one after another on the system, with outcomes $A_i$:
$\Pr(A_{<i}=\mathbf a,A_i=k)=\nrm{P_i^kQ_{\mathbf a}\xi}_F^2=:p(\mathbf a,k)$.
\begin{enumerate}[label=\arabic*.,itemsep=1pt]
\item The reference is untouched, so $S(\rho)\ge I(A_1\cdots A_N:\mathrm{Ref})=\sum_iI(A_i:\mathrm{Ref}\mid A_{<i})$.
\item By $\nrm{(c-I)Y}_F^2=3\sum_{k\ne0}\nrm{P^kY}_F^2$ and $[c_i,Q_{\mathbf a}]=0$, $\Pr(A_i\ne0)=\frac13\sum_{\mathbf a}\nrm{Q_{\mathbf a}(c_i-I)\xi}_F^2=\Delta_i^2/3$, an equality.
\item $x_iP_i^{-k}x_i=P_i^k$, $x_i$ commutes with $Q_{\mathbf a}$, and right multiplication by $x_i$ is unitary, so $\sqrt{p(\mathbf a,-k)}=\nrm{P_i^kQ_{\mathbf a}x_i\xi}_F$ and
$\sqrt{p(\mathbf a,k)}=\nrm{P_i^kQ_{\mathbf a}\xi x_i}_F$; hence $\sum_{\mathbf a,k}(\sqrt{p(\mathbf a,k)}-\sqrt{p(\mathbf a,-k)})^2\le\sum_{\mathbf a}\nrm{Q_{\mathbf a}[x_i,\xi]}_F^2=\eta_i^2$, and the terms
$k=1$ and $k=2$ are equal. The inequality $H(r_0,r_1,r_2)\ge s-(\sqrt{r_1}-\sqrt{r_2})^2/\ln2$, $s=r_1+r_2$, of~\cite[equation~(16)]{DouglasMghirbiC} is homogeneous, so it applies to
the unnormalized weights given each $\mathbf a$, and $H(A_i\mid A_{<i})\ge\Delta_i^2/3-\eta_i^2/(2\ln2)$.
\item Let the reference measure $c_i$ from the right, with outcome $B_i$. Left and right multiplications commute, and by $\nrm{[c,Y]}_F^2=3\sum_{k\ne l}\nrm{P^kYP^l}_F^2$ and
$[c_i,Q_{\mathbf a}\xi]=Q_{\mathbf a}[c_i,\xi]$, $\Pr(A_i\ne B_i)=\theta_i^2/3$, an equality.
\item By data processing and Fano's inequality for three values, with the concavity of $h_2$ over $\mathbf a$,
$I(A_i:\mathrm{Ref}\mid A_{<i})\ge I(A_i:B_i\mid A_{<i})\ge H(A_i\mid A_{<i})-h_2(\theta_i^2/3)-\theta_i^2/3$.
\end{enumerate}
Summing over $i$ proves the bound. On the flat state of a tensor product of standard representations, $\Delta_i^2=3$ and $\theta_i=\eta_i=0$, so the bound is $N=S$.
\end{proof}

\begin{proof}[Proof of Lemma~\ref{lem:centrality}]
$U$ agrees with $U_\pi$ on $\C^d\otimes K$ and distinct slots occupy distinct matrix positions, so $T(\rho)=\sum_a(|c_a|^2/d)\rho_a$ with $\rho_a=\pi(w_a)\rho\pi(w_a)^*$, the weights
summing to $1$. Every $\rho_a$ has the spectrum of $\rho$, so $a=S(\bar\rho)-\sum_a(|c_a|^2/d)S(\rho_a)=\sum_a(|c_a|^2/d)D(\rho_a\|\bar\rho)$ with $\bar\rho=T(\rho)$ (Holevo's identity). The
anchor slot of a symbol $w$, and that of the empty word, have coefficient $1$ and hence weight $1/d$, so $D(\rho_w\|\bar\rho)\le d\,a$ and $D(\rho\|\bar\rho)\le d\,a$. For states,
$\nrm{\sqrt\sigma-\sqrt\tau}_F^2=2-2\Tr\sqrt\sigma\sqrt\tau\le-2\ln\Tr(\sqrt\sigma\sqrt\tau)\le D(\sigma\|\tau)$ in nats, since the Petz--R\'enyi divergences increase with the order. So
$\nrm{\sqrt{\rho_w}-\sqrt{\bar\rho}}_F$ and $\nrm{\sqrt\rho-\sqrt{\bar\rho}}_F$ are at most $\sqrt{d\,a\ln2}$, and since $\pi(w)\xi\pi(w)^*=\sqrt{\rho_w}$ and right multiplication by a unitary preserves
the norm, $\chi(\pi(w))=\nrm{\sqrt{\rho_w}-\sqrt\rho}_F\le2\sqrt{d\,a\ln2}$.
\end{proof}

\begin{theorem}[Explicit genuine-round constants]\label{thm:genuine-constants}
Let $(U,\rho)$ be a representation round of $\Phi_H$ with distance at most $\frac18$ and entropy production $a$, with any leakage, and let $N_{\rm ex}$ be
the largest $N$ for which $(c_i,x_i)_{i<N}$ satisfy the hypotheses of Lemma~\ref{lem:sector}. Then
\[
S(\rho)\ge0.375\min\Bigl(N_{\rm ex},\ \frac1{60\delta_{\rm cent}}-\frac23\Bigr),\qquad \delta_{\rm cent}=\max\bigl(\chi(\pi(\sfa)),\chi(\pi(\alpha))\bigr)\le16.981\sqrt a .
\]
For a genuine round $N_{\rm ex}=\infty$, so
\[
S(\rho)\ge\frac{a^{-1/2}}{2717}-\frac14\qquad\text{and}\qquad S+n\,a\ge0.0097\,n^{1/3}-\frac14\quad\text{for every }n\ge1 .
\]
\end{theorem}

\begin{proof}[Proof of Theorems~\ref{thm:genuine} and~\ref{thm:genuine-constants}]
\emph{Displacement.} $\gamma$ is a symbol, a prefix of $(\alpha\sfa\alpha\sfa^{-1})^3$, with an anchor $z_\gamma$, and the empty word has an anchor $0$. By (F1),
$\bra{z_\gamma}\Psi(\ket{z_\gamma}\bra0)\ket0=\Tr(\rho\pi(\gamma))$, while for $\Phi_H$ it is $\tau(\lambda(\gamma))=0$, since $\gamma\ne1$ in $H_3$. By (F2), $|\Tr\rho\pi(\gamma)|\le2u\le\frac14$, so
$E(\pi(\gamma))^2=2-2\mathrm{Re}\Tr(\rho\pi(\gamma))\ge\frac32$.

\emph{Transport.} $c_i=\pi(\sfa)^{3i}\pi(\gamma)\pi(\sfa)^{-3i}$, so by (W4) and (W2), $\Delta_i\ge\sqrt{1.5}-2\chi(\pi(\sfa)^{3i})\ge\sqrt{1.5}-6i\delta_{\rm cent}$. The words of $c_i$ and $x_i$ have lengths
$6i+4$ and $6i+1$ in $\sfa^{\pm1}$ and $\alpha$, so by (W2) $\theta_i\le(6i+4)\delta_{\rm cent}$ and $\eta_i\le(6i+1)\delta_{\rm cent}$.

\emph{Copies.} Let $N=\min(N_{\rm ex},\lfloor(0.1/\delta_{\rm cent}+2)/6\rfloor)$. For $i<N$, $\theta_i,\eta_i\le(6N-2)\delta_{\rm cent}\le0.1$ and $6i\delta_{\rm cent}\le0.1$. Lemma~\ref{lem:sector},
which holds on any Hilbert space for states of finite rank, applies directly with these $\Delta_i,\theta_i,\eta_i$. Each term is at least $(\sqrt{1.5}-0.1)^2/3-0.01/(2\ln2)-h_2(\frac1{300})-\frac1{300}=0.37891>0.375$, so
$S(\rho)\ge0.375N\ge0.375\min(N_{\rm ex},1/(60\delta_{\rm cent})-\frac23)$.

\emph{Centrality.} $\sfa$ and $\alpha$ are symbols, so Lemma~\ref{lem:centrality} gives $\delta_{\rm cent}\le2\sqrt{104\ln2}\,\sqrt a=16.981\sqrt a$.

\emph{Genuine rounds.} In $H_3$, $\sfa^{3i}\gamma\sfa^{-3i}$ is a three-cycle and $\sfa^{3i}\alpha\sfa^{-3i}$ a transposition on a triple of consecutive points of the line through the
first two rays, shifted by $3i$, with $\alpha\gamma\alpha=\gamma^{-1}$~\cite[Section~4]{DouglasMghirbiB}; triples for different $i$ are disjoint. So the relations of
Lemma~\ref{lem:sector} hold in every representation of $H_3$, $N_{\rm ex}=\infty$, and $S\ge0.375(a^{-1/2}/(60\cdot16.981)-\frac23)\ge a^{-1/2}/2717-\frac14$. Finally
$\min_{x>0}(Ax^{-1/2}+nx)=(2^{1/3}+2^{-2/3})A^{2/3}n^{1/3}$, which is at least $0.0097\,n^{1/3}$ for $A=1/2717$.
\end{proof}

\begin{theorem}[Explicit genuine-service constants]\label{thm:genuine-services-constants}
For $n\ge n_0$, every genuine service of $\Phi_H$ with error
$\epsilon\le1/16$ and $J\le n\log26/2$ has $Q\ge n^{1/3}/385$
if its discard is zero, and $Q\ge n^{1/3}/420$ if its discard is
at most $1/n^2$. Consequently the lower constant in
Theorem~\ref{thm:houghton-genuine} may be taken as $1/420$.
The nongenuine-round threshold in Theorem~\ref{thm:houghton-reduction}(c)
may be taken as $S+na<0.0097n^{1/3}-1/4$.
\end{theorem}

\begin{proof}[Proof of Theorems~\ref{thm:genuine-services} and~\ref{thm:genuine-services-constants}]
If $U$ is genuine on $V$ and $\rho$ is supported in $V$, the round $(U,\rho)$ has the channel of the representation round $(U_\pi,J\rho J^*)$, since the partial trace over the memory
ignores the isometry $R$; $T_U(\rho)=RT_\pi(J\rho J^*)R^*$ has the spectrum of $T_\pi(J\rho J^*)$, and $S(J\rho J^*)=S(\rho)$. So Theorem~\ref{thm:genuine-constants} applies to $(U,\rho)$.

\emph{Discard $0$.} By Theorem~\ref{thm:lower-constants} some round $t$ of the service has distance at most $\frac1{16}$, $S(\rho_t)\le Q$ and $n\,a_t\le6.4Q+13.52+0.134\log n$, and the
conditioned memory state $\rho_t$ is supported in $V_t$, since conditioning only shrinks supports. Theorem~\ref{thm:genuine-constants} gives $Q+\frac14\ge a_t^{-1/2}/2717$. If
$Q\ge14+0.14\log n$, then $n\,a_t\le7.4Q$, so $(Q+\frac14)Q^{1/2}\ge n^{1/2}/(2717\sqrt{7.4})\ge n^{1/2}/7392$, and with $Q+\frac14\le(1+\frac1{56})Q$, $Q\ge n^{1/3}/384$. If
$Q<14+0.14\log n$, then $a_t=O(\log n/n)$ and Theorem~\ref{thm:genuine-constants} gives $Q+\frac14\ge c(n/\log n)^{1/2}$, which is impossible for $n\ge n_0$.

\emph{Discard at most $1/n^2$.} Let $t$ be the round of Theorem~\ref{thm:lower-constants} and $\rho=\rho_t$. The test observer's memory state at the start of round $t$, conditioned on
passing, an event of probability at least $1-\epsilon$, puts mass $\delta'\le\delta$ outside $V:=V_t$. Let $P$ project onto $V$ and $\hat\rho=P\rho P/\Tr(P\rho)$. If $\psi$ purifies $\rho$, then
$(P\otimes I)\psi/\sqrt{1-\delta'}$ purifies $\hat\rho$ and has overlap $\sqrt{1-\delta'}$ with $\psi$, so $\nrm{\rho-\hat\rho}_1\le2\sqrt{\delta'}\le2\sqrt\delta$, by the pure-state trace-distance formula and contraction under partial trace. The round $(U,\hat\rho)$ is genuine, and:
its distance is at most $\epsilon+\frac12\nrm{\rho-\hat\rho}_1\le\frac1{16}+\frac1n\le\frac18$ for $n\ge16$, since the channel is linear in the bath; $T$ is a channel on the $Q$-qubit memory,
so $\frac12\nrm{T(\rho)-T(\hat\rho)}_1\le\sqrt\delta\le\frac12$, and Audenaert's sharp form of Fannes' inequality~\cite[Theorem~1]{Audenaert2007} gives
$|S(T(\rho))-S(T(\hat\rho))|\le\sqrt\delta\,Q+h_2(\sqrt\delta)$, and likewise for $S(\rho)-S(\hat\rho)$, so $a(\hat\rho)\le a(\rho)+2\sqrt\delta\,Q+2h_2(\sqrt\delta)$; and $S(\hat\rho)\le Q$. With
$\delta\le1/n^2$ and $n\,h_2(1/n)\le\log n+\log e$, $n\,a(\hat\rho)\le8.4Q+16.41+2.134\log n$. If $Q\ge16.41+2.134\log n$ then $n\,a(\hat\rho)\le9.4Q$, and
Theorem~\ref{thm:genuine-constants} gives $(Q+\frac14)Q^{1/2}\ge n^{1/2}/(2717\sqrt{9.4})\ge n^{1/2}/8331$, so $Q\ge n^{1/3}/415.2\ge n^{1/3}/420$. Otherwise $a(\hat\rho)=O(\log n/n)$ and
Theorem~\ref{thm:genuine-constants} gives a contradiction for $n\ge n_0$.
\end{proof}

With discard of order $\epsilon^2/n$ the term $2\sqrt\delta\,Q$ is of order $\epsilon Q/\sqrt n$, which dominates the budget $Q/n$, and the same argument gives only $Q\gtrsim n^{1/6}$.

\begin{proof}[Proof of Theorem~\ref{thm:houghton-genuine}]
The lower bound is Theorem~\ref{thm:genuine-services}. For the upper bound take the service of Theorem~\ref{thm:houghton-groups} with $q=3$; for $n\ge n_0(\epsilon)$ its memory
is $O(n^{1/3}(\log n)^{2/3}(\log(n/\epsilon))^{1/3})$. Its round unitary is $U_t=R''W_cR'$, where $R'$ (class permutation and twirls) and $R''$ (re-fold with its completion, twirls)
act on the memory alone and $W_c$ is the round. $R'$ maps the window subspace $V$ (window units, slots and ancilla) onto itself, and on $\C^d\otimes V$ the round is the gadget unitary
of the right regular representation $\pi$ of $H_3$, tensored with the identity on slots and ancilla. So $U_t$ is genuine on $V$, with $J$ the identification of $V$ with a
subspace of $\ell^2(H_3)\otimes\text{slots}\otimes\text{ancilla}$ composed with $R'$, and $R=R''$ on the images of window points under slot words, which lie in $\Omega_W$. By
Theorem~\ref{thm:fibred}, for every observer the memory has mass at most $3/(4n^2)\le(1-\epsilon)/n^2$ outside $V$ at the start of every round. So the service is genuine with
discard at most $1/n^2$.
\end{proof}

\begin{proof}[Proof of Corollary~\ref{cor:growth-route}]
The labels of $\Phi_H$ generate $H_3$, which is elementary amenable. It is not residually finite, since every finite-dimensional representation kills its three-cycles
(Remark~\ref{rem:houghton-map}), so it is not virtually nilpotent, and by Chou's theorem~\cite[Theorem~3.2]{Chou1980} it has exponential growth. The regular action is free, so the
words $w_x$ of Proposition~\ref{prop:growth} give distinct points, and the proposition applies to every finitely supported bath.
\end{proof}

\begin{proof}[Proof of Proposition~\ref{prop:genuine-vacuous}]
By Theorem~\ref{thm:genuine-constants}, a genuine round at distance at most $\frac18$ has $1+S+n\,a\ge0.0097\,n^{1/3}+\frac34$. For $n\ge n_1$ the service of
Theorem~\ref{thm:houghton-groups} with $q=3$ has memory $O(n^{1/3}(\log n)^{2/3}L^{1/3}+L\log L)+O(\log(n/\epsilon))$ with $L=O(\log(n/\epsilon))$, which is at most
$K\log^2(n/\epsilon)\,n^{1/3}$, hence at most $K'\log^2(n/\epsilon)(1+S+n\,a)$. For $2\le n<n_1$, Proposition~\ref{prop:ceiling} gives $\Qs(n,\epsilon)\le n_1\lceil\log26\rceil$.
\end{proof}

\begin{proof}[Proof of Proposition~\ref{prop:ceiling-houghton}]
By~\cite[Proposition~4.14]{DouglasMghirbiB}, for every $u_0>0$ there is $C(u_0)$ with $\log D^{\rm ent}_{\Phi_H}(u_0,v)\le C(u_0)v^{-1/3}$ for $0<v\le1$, where $D^{\rm ent}$ is the least bath
dimension of a round with distance at most $u_0$ and $\ell+a\le v$; the rounds of its proof have flat baths, so $a=0$ and $\ell\le v$, and $S$ is at most the logarithm of the bath
dimension. With $u_0=\epsilon$ and $v=\frac{32}{15}\epsilon/n$ they are admissible in $L_{\Phi_H}(n,\epsilon)$, so $L_{\Phi_H}(n,\epsilon)\le C(\epsilon)(15n/(32\epsilon))^{1/3}$, and with
$u_0=\frac1{16}$, $v=\frac2{15n}$ the same gives $\Ls_{\Phi_H}(n)\le C\,n^{1/3}$.
\end{proof}

\paragraph{The single budget (Remark~\ref{rem:quarter}).} At the rank rate, \cite[Theorem~5.2]{DouglasMghirbiB} with Proposition~\ref{prop:purity} reads a round with
$\ell+a\le t_n\le\frac{16}{15}(3Q+8+\log n)/n$. A lower bound $\Qs\ge Q_0$ by this method requires that the least bath entropy of a round with distance $\epsilon/(1-\epsilon)$ and
$\ell+a\le t_n(Q)$ exceed $Q$ for all $Q<Q_0$. By the proof of Proposition~\ref{prop:ceiling-houghton} that bath entropy is at most $C\,t_n(Q)^{-1/3}$, so the requirement fails as
soon as $C(\frac{16}{15}(3Q+8+\log n)/n)^{-1/3}\le Q$, which holds for $Q\ge C'n^{1/4}$.

\begin{proof}[Proof of Theorem~\ref{thm:houghton-reduction}]
(a) Let $\Ls_{\Phi_H}(n)\ge n^{1/3}/p(n)$. As in the proof of Proposition~\ref{prop:genuine-vacuous}, $\Qs_{\Phi_H}(n,\epsilon)\le K\log^2(n/\epsilon)\,n^{1/3}$ for all $n\ge2$ (for small $n$
by Proposition~\ref{prop:ceiling}), so $\Qs_{\Phi_H}(n,\epsilon)\le Kp(n)\log^2(n/\epsilon)(1+\Ls_{\Phi_H}(n))$, the upper half of the law; the lower half is Theorem~\ref{thm:lower-constants},
and it gives $\Qs_{\Phi_H}(n,\epsilon)\ge(n^{1/3}/p(n)-13.6-0.14\log n)/7.4$. So $\Qs_{\Phi_H}(n,\epsilon)=n^{1/3}\cdot\mathrm{polylog}(n/\epsilon)$.

(b) If the law holds, $\Ls_{\Phi_H}(n)\ge\Qs_{\Phi_H}(n,\epsilon)/(K(\log(n/\epsilon))^\kappa)-1$, and under Conjecture~\ref{conj:houghton} the right side is $n^{1/3-o(1)}$.

(c) This is Theorem~\ref{thm:genuine-constants}.
\end{proof}

\subsection{Constant-field geometry and sector-isotropic sources}

We give a direct proof of Corollary~\ref{thm:number-sectors}. The geometric input is
the existing construction of~\cite[Section~6.2]{DouglasMghirbiA}:
$\gamma$ is exactly a three-cycle at every clock; the $r_4$ holonomy
is a product of disjoint rank-one phase rotations with angles
$F_R=s_R\pi/|R|$. There are $M$ regions, each of area between
$m_M=1+(M-1)(M-2)/2$ and $M^2$. These are model identities,
not conclusions for an arbitrary round.

\paragraph{Sector characters force entropy.}
Initially suppose the source is clock diagonal. Set
$v_j=j(L-j)/[L(L-1)]$ and $V_{\mathbf j}=\sum_l v_{j_l}$.
The normalized three-cycle character in degree $j$ is
\[
 \chi_j=1-3v_j.
\]
Indeed its eigenvalues are $1$ repeated $L-2$ times and
$\omega,\omega^2$; the generating polynomial is
$(1+t)^{L-2}(1-t+t^2)=(1+t)^L-3t(1+t)^{L-2}$.
These characters are nonnegative for $L=4M+4$. Conditional sector
flatness removes conjugations of the cycle and gives
$\chi_\gamma=\E_p\prod_l(1-3v_{j_l})$.
The identity and $\gamma$ singleton anchors imply
$0\leq\chi_\gamma\leq2u\leq1/8$.
In the rest of the lower proof it suffices to assume
$\chi_\gamma\leq1/4$. The product is at least $1-3V$, so
$\E V\geq1/4$; on $V<1/12$ it is greater than $3/4$,
giving $\Pr(V\geq1/12)\geq2/3$.
For $m=\min(j,L-j)$, $\binom Lj\geq2^m$ and $Lv_j\leq m$.
The first inequality follows termwise from
$(L-i)/(m-i)\geq2$ for $0\leq i<m$.
Consequently, with $D_{\mathbf j}=\prod_l\binom L{j_l}$,
\begin{equation}\label{eq:sector-entropy-floor}
 S=H(p)+\E\log_2D_{\mathbf j}\geq L\E V\geq L/4=M+1.
\end{equation}
Source entropy, rather than logarithmic source rank, pays this bound.

\paragraph{Clock bias costs actual production.}
Every literal slot translates the clock and applies a degree-preserving
one-particle unitary. It preserves each $\bigotimes_l\tau_{j_l}$.
The actual Choi bath map is therefore classical clock convolution,
$Kp=\sum_s w_sp_s$, leaving degrees fixed. Identity and both
coordinate steps have weights at least $1/d$ from their singleton
slots, with $d=104$. Thus
\[
 a=H(Kp)-H(p)=\sum_s w_sD_2(p_s\Vert Kp).
\]
The unchanged average $\log D_{\mathbf j}$ cancels exactly.
Root relative entropy and the triangle inequality bound each coordinate
root displacement by $4d(\ln2)a$. The $M$-torus Fourier Poincar\'e
constant is $4\sin^2(\pi/M)\geq16/M^2$. Applied separately to
each degree block and then summed, it gives
\[
 \sum_{c,\mathbf j}(\sqrt{p(c,\mathbf j)}
       -\operatorname{mean}_c\sqrt{p(c,\mathbf j)})^2
       \leq(d\ln2)M^2a/2.
\]
Put $w(\mathbf j)=\sum_cp(c,\mathbf j)$ and
$q(c,\mathbf j)=w(\mathbf j)/M^2$.
In a nonzero degree block let
$x=M\operatorname{mean}_c\sqrt p/\sqrt w\in[0,1]$.
Its variance is $w(1-x^2)$, whereas its squared root distance
from the uniform clock law is $2w(1-x)\leq2w(1-x^2)$.
Cauchy--Schwarz therefore proves
\begin{equation}\label{eq:joint-clock-degree}
 \operatorname{TV}(p,q)\leq
 \nrm{\sqrt p-\sqrt q}_F\leq M\sqrt{d(\ln2)a}.
\end{equation}
This argument retains all clock--degree and flavour correlations.

\paragraph{A protected region forces leakage.}
Let $B_0=\{0,\ldots,\lfloor M/4\rfloor\}^2$.
It has at least $M^2/16$ clocks and is contained in all $M$ regions.
To check containment, the $M-1$ bulk regions indexed by
$5\leq t\leq M+3$ have source $(M-1,M-1)$, the column
$(M-1,j)$ for $1\leq j\leq M-2$, and bulk points
$0\leq i\leq M-3$, $0\leq j\leq M-2$, $i+j\leq t+M-8$.
Their areas are $(M-1)^2-h(h+1)/2$, $h=M+3-t$.
The last region contains $i,j\geq0$, $i+j\leq M-3$, as well as
$(0,M-2)$. The inequality $2\lfloor M/4\rfloor\leq M-3$
proves the claim. At a clock in $B_0$,
$\sum_RF_R^2\geq\pi^2/M^3$.
There are at most $M<L/4$ nonzero phase modes and
$\sum_R|F_R|\leq M\pi/m_M\leq\pi$.

For a uniformly sampled degree-$j$ subset of the $L$ modes set
$X=\sum_RF_RI_R$. Sampling without replacement gives
\[
 \operatorname{Var}X=v_j\left[\sum_RF_R^2-
             (\sum_RF_R)^2/L\right]\geq\tfrac34v_j\sum_RF_R^2.
\]
For independent copies, $|X-X'|\leq\pi$ and
$1-|\E e^{iX}|^2=\E[1-\cos(X-X')]
\geq4\operatorname{Var}X/\pi^2$.
Hence $|\E e^{iX}|\leq
\exp[-3v_j\sum_RF_R^2/(2\pi^2)]$.
Multiplication across conditionally flat flavours shows that on
$B_0\cap\{V\geq1/12\}$ the sector relator squared defect obeys
\[
 e_{c,\mathbf j}^2\geq2(1-e^{-1/(8M^3)})\geq1/(8M^3).
\]
No bound on the total phase across flavours is needed.

Project the physical round onto the canonical coefficients and
let $\Lambda$ be its bath leakage operator as in
Proposition~\ref{prop:kernel-concentration}. Its identity singleton
and closing $r_4$ triangle have the same target label and squared
weights $1,1/2$. A common projected endpoint $z$ must therefore pay
\[
 \inf_z\{\|\psi-z\|^2+\tfrac12\|\pi(r_4)\psi-z\|^2\}
       =\tfrac13\|(\pi(r_4)-I)\psi\|^2.
\]
Every other slot contributes nonnegative error, so
$\Lambda\geq(\pi(r_4)-I)^*(\pi(r_4)-I)/3$.
Leakage on the protected event is at least $1/(24dM^3)$.
Its $q$-mass is at least $(1/16)(2/3)=1/24$.
Set $A=1/(2304d\ln2)$, $B=1/(1152d)$.
Equation~\eqref{eq:joint-clock-degree} gives
\begin{equation}\label{eq:sector-flux-dichotomy}
 a\leq A/M^2\quad\Longrightarrow\quad\ell\geq B/M^3,
\end{equation}
because its total-variation error is at most $1/48$.
At leakage $\ell\leq c/n$, either $M\geq(Bn/c)^{1/3}$ or
$S+na\geq M+nA/M^2\geq3A^{1/3}n^{1/3}/2^{2/3}$.
This proves the diagonal-source lower with constant
$\min\{(B/c)^{1/3},3A^{1/3}/2^{2/3}\}$.

\paragraph{Coherent clocks and flags.}
Clock pinching commutes with the physical bath map: each slot has
one fixed clock translation. Data processing of
$D_2(\rho\Vert\Delta_c\rho)=S(\Delta_c\rho)-S(\rho)$ gives
$a(\Delta_c\rho)\leq a(\rho)$.
Every literal slot assigned a target label has that label's clock
translation. Its projected coefficient and residual have the same
translation; each residual square preserves the clock. Thus $\Lambda$
is clock block diagonal, and pinching preserves leakage.
The $\gamma$ character has zero translation and is preserved as well.
To retain the entropy floor directly, coherently copy the clock to an
analytical register. Its marginals have entropies
$S(\Delta_c\rho)$ and $H(p_c)$, so Araki--Lieb gives
$S(\rho)\geq\sum_cp_cS(\rho_c)$, where $\rho_c$ is the normalized
clock-diagonal block. Each $\rho_c$ is an orthogonal degree-sector
mixture, hence
\[
 S(\rho)\geq\sum_{c,\mathbf j}p(c,\mathbf j)\log_2D_{\mathbf j}
        \geq L\E V\geq L/4=M+1.
\]
There is no clock-entropy subtraction or free pinching assumption.
If $a(\rho)\leq A/M^2$, the pinched source satisfies
\eqref{eq:sector-flux-dichotomy}; otherwise optimize $M+nA/M^2$.
This proves the coherent-source lower with the same constant as the
clock-diagonal case.

For a finite block-diagonal round with coherent flag source, let
$b_i=\Tr P_i\rho$ and $\rho_i=P_i\rho P_i/b_i$.
Channel and leakage depend only on diagonal flag blocks.
All $\chi_i\geq0$, and the overall anchor test gives
$\sum_ib_i\chi_i\leq1/8$. Components with $\chi_i\leq1/4$
have mass at least $1/2$; those with $\ell_i\leq8/(15n)$ have
mass at least $3/4$ by Markov. Their intersection has mass at least
$1/4$. Flag pinching commutes with the bath map, so data processing
gives $a(\rho)\geq\sum_ib_i a(\rho_i)$.
Coherently copying the flag is an analytical isometry. Its marginals
have entropies $H(b)+\sum_ib_iS(\rho_i)$ and $H(b)$;
Araki--Lieb gives $S(\rho)\geq\sum_ib_iS(\rho_i)$.
Apply the coherent clock bound to the good intersection. These
comparisons charge the original source; no physical pinching,
purifier or flag supply is introduced.

\paragraph{Replacement and upper bound.}
The triangle-good corner families of~\cite[Proposition~4.10]{DouglasMghirbiB} have
source entropy $O(m)$, room $O(m^{-2})$ and scalar gap at most $1/2$.
Their actual gain is bounded by that room. Fixed untwisted tensor
amplification gives target error at most $1/16$, zero leakage and
cost $O(m+nm^{-2})$. Choose $m\asymp n^{1/3}$; small $n$ use the
pure exact baseline. The source lies entirely in the new round's
exact leak kernel. This different round need not stay close to the
smoothed source or belong to its geometry.
Finally the fully flat smoothed models themselves, with fixed flavour
count selected for coarse error $1/16$, have $a=0$, leakage
$O(M^{-3})$ and entropy $O(M)$
\cite[Proposition~4.14]{DouglasMghirbiB}.
Taking $M\asymp n^{1/3}$ proves the class's matching static upper.
Neither lower is an all-family operational memory statement.

\subsection{Collective orbital invariance}

We prove Theorem~\ref{thm:collective-sources}, retaining precisely the
geometry above but dropping product-sector flatness. For an invariant
internal state $\eta$, let $Q$ project onto the nontrivial part of the
restriction to $SU(2)$ acting on a two-dimensional one-particle plane,
and put $p_\eta=\Tr\eta Q$. All planes are conjugate in $U(L)$,
so this number does not depend on the plane.

\paragraph{Entropy without flattening multiplicities.}
The three-cycle is a product of two $SU(2)$ quarter turns on planes
$(1,2)$ and $(2,3)$: each sends its first basis vector to the second
and the second to minus the first. The product sends $e_1$ to $e_2$,
$e_2$ to $e_3$ and $e_3$ to $e_1$.
For either turn, $(\Pi(A)-I)^*(\Pi(A)-I)\leq4Q$.
The Hilbert--Schmidt triangle inequality and invariance give
\[
 2(1-\Re\Tr\eta\Pi(\gamma))
 =\nrm{(\Pi(\gamma)-I)\sqrt\eta}_F^2\leq16p_\eta.
\]
Choose $t=\lfloor L/2\rfloor$ disjoint planes. The restriction to
their commuting $SU(2)$ groups decomposes into tensor products of
irreducibles, with arbitrary multiplicity spaces. Schur's lemma
makes an invariant state flat only on the irreducible factors:
\[
 \eta=\bigoplus_{\mathbf m}w_{\mathbf m}
 \left[\bigotimes_{i=1}^t\frac{I_{m_i+1}}{m_i+1}
                  \otimes\sigma_{\mathbf m}\right].
\]
Here the multiplicity states $\sigma_{\mathbf m}$ need not be flat
or separable. Every nontrivial factor costs at least one bit, hence
$S(\eta)\geq\sum_i\Pr(m_i>0)=tp_\eta$.
For completeness, the finite-dimensional unitary $SU(2)$ irreducibles
have weights $m,m-2,\ldots,-m$ and dimension $m+1$.
With Lie algebra generators $E,F,H$, a maximal-weight vector obeys
$EF^jv=j(m-j+1)F^{j-1}v$. Positivity of the unitary norm and termination
force an integer $m\geq0$ and a chain $j=0,\ldots,m$; invariant
orthogonal complements complete the decomposition.

Set $\bar\eta=\sum_cq_c\rho_c$ and
$\chi=\Tr\bar\eta\Pi(\gamma)$. The clock blocks of $\gamma$ are
conjugate three-cycles, so invariance makes this their common trace
test. It is real, and the anchors give $|\chi|\leq2u\leq1/8$.
The direct conditional-entropy argument above gives
\[
 S(\rho)\geq\sum_cq_cS(\rho_c)
       \geq t p_{\bar\eta}.
\]
Even the weaker $\chi\leq1/4$ therefore implies
$p_{\bar\eta}\geq3/32$ and $S(\rho)\geq bM$, $b=3/16$.

\paragraph{An operator-valued clock law.}
Put $\bar\rho=\Delta_c\rho=\bigoplus_cR_c$,
$R_c=q_c\rho_c$. Invariance makes every internal literal unitary
act trivially on these blocks; the bath map is clock convolution
$K\bar\rho=\sum_sw_s\bar\rho_s$.
The translated states have the same spectrum, so the Holevo identity gives
\[
 \bar a=S(K\bar\rho)-S(\bar\rho)
          =\sum_sw_sD_2(\bar\rho_s\Vert K\bar\rho).
\]
The inequality $\nrm{\sqrt X-\sqrt Y}_F^2\leq D_{\rm nat}(X\Vert Y)$
follows by diagonalizing both states: their overlap-weighted eigenvalue
arrays give exactly the quantum relative entropy and affinity,
and Jensen gives $D_{\rm nat}\geq-2\ln\Tr\sqrt X\sqrt Y$
$\geq2(1-\Tr\sqrt X\sqrt Y)$.
Identity and coordinate weights at least $1/d$ bound each root
shift energy by $4d(\ln2)\bar a$.
Apply the torus Poincar\'e inequality to the Hilbert-space field
$X_c=\sqrt{R_c}$, not to simultaneously diagonalized densities.
With $A_0=M^{-1}\sum_cX_c$,
\[
 V=\sum_c\nrm{X_c-M^{-2}\sum_bX_b}_F^2
   =1-\Tr A_0^2\leq\tfrac12(d\ln2)M^2\bar a.
\]
Operator Jensen gives $A_0^2\leq\sum_cR_c=\bar\eta$;
square-root monotonicity gives $A_0\leq\sqrt{\bar\eta}$.
Thus $\Tr A_0\sqrt{\bar\eta}\geq\Tr A_0^2$, without commutation.
For $q=(I/M^2)\otimes\bar\eta$ the root distance is at most $2V$.
The square-root trace-distance inequality yields
\begin{equation}\label{eq:operator-clock-law}
 \tfrac12\nrm{\bar\rho-q}_1\leq M\sqrt{d(\ln2)\bar a}
                         \leq M\sqrt{d(\ln2)a}.
\end{equation}
The final inequality, and preservation of leakage and $\chi$,
follow from clock-pinching covariance and data processing exactly
as above. No internal twirl or commuting-density assumption is used.

\paragraph{A spin-uniform flux estimate.}
For a nontrivial $SU(2)$ irrep, $m\geq1$, and $|\theta|\leq\pi/4$,
its normalized character satisfies
\begin{equation}\label{eq:uniform-spin-character}
 \left|\frac{\sin((m+1)\theta)}{(m+1)\sin\theta}\right|
          \leq e^{-\theta^2/32}.
\end{equation}
To verify this uniformly, set $N=m+1$. If $N|\theta|\leq1$,
the character is the mean of $e^{iX}$ for weights
$X=(N-1-2j)\theta$. Its variance is
$(N^2-1)\theta^2/3\geq\theta^2$, and $|X-X'|\leq2<\pi$.
The cosine inequality gives the stronger modulus bound
$e^{-2\theta^2/\pi^2}$. If $N|\theta|>1$, then
$|\sin x|/x\leq\sin1$ for $x\geq1$, by monotonicity up to $\pi$
and $1/x$ thereafter. Also $\theta/\sin\theta<9/8$ and
$\sin1\leq101/120$. Hence the modulus is at most $909/960$,
smaller than $e^{-\theta^2/32}\geq1-\pi^2/512$.
The value at zero is understood by continuity.

Let a one-particle holonomy $W$ have angles $F_j$ on at most $M$
disjoint phase modes. Pair each with an unused mode, and let $C$ swap
the partners. Then $h=WCW^*C^*$ is a product of disjoint $SU(2)$
rotations of angles $F_j$. Restriction of $\eta$ to their product gives
a weighted sum of products of characters. Using
\eqref{eq:uniform-spin-character}, the triangle inequality and
$2(1-e^{-x})\geq x$ for $0\leq x\leq1$, when $\sum_jF_j^2\leq1$,
\[
 \nrm{(\Pi(h)-I)\sqrt\eta}_F^2
        \geq p_\eta\sum_jF_j^2/32.
\]
Invariance under $C$ and $W$ bounds the left root norm by
$2\nrm{(\Pi(W)-I)\sqrt\eta}_F$. Consequently
\begin{equation}\label{eq:invariant-flux-floor}
 \nrm{(\Pi(W)-I)\sqrt\eta}_F^2
        \geq p_\eta\sum_jF_j^2/128.
\end{equation}
This estimate is independent of all spin, flavour and multiplicity
dimensions; it does not assume nonnegative three-cycle characters.

\paragraph{Optimization and coherent flags.}
The existing fields have $\sum_jF_j^2<1$ and $|F_j|<\pi/4$
for $M\geq12$. On $B_0$, $\sum_jF_j^2\geq\pi^2/M^3$.
The same two-slot leakage lower and~\eqref{eq:invariant-flux-floor}
give, for any fixed plane projection $Q$,
\[
 \ell\geq\frac{\pi^2}{384dM^3}
                 \Tr\bar\rho(P_{B_0}\otimes Q).
\]
Invariance makes $Q$'s conditional expectation plane-independent.
For $q$ in~\eqref{eq:operator-clock-law}, this probability is at least
$p_{\bar\eta}/16\geq3/512$.
Set $A=9/(1048576d\ln2)$ and $B=\pi^2/(131072d)$.
If $a\leq A/M^2$, the comparison changes it by at most $3/1024$,
so $\ell\geq B/M^3$.
With $\ell\leq c/n$ and $S\geq bM$, optimization gives
\[
 S+na\geq c_H(c)n^{1/3},\qquad
 c_H(c)=\min\left\{b(B/c)^{1/3},
                   \frac{3(b^2A)^{1/3}}{2^{2/3}}\right\}.
\]

For coherent model flags the conditional entropy/gain inequalities
above still hold. Here each $\chi_i\in[-1/2,1]$, because the
three-cycle has order three and is conjugate to its inverse.
The overall anchor bound $\sum_ib_i\chi_i\leq1/8$ gives mass
at least $1/6$, not $1/2$, on $\chi_i\leq1/4$.
Markov gives mass at least $11/12$ on $\ell_i\leq8/(5n)$;
the intersection has mass at least $1/12$.
The lower is therefore $c_H(8/5)n^{1/3}/12$.
The fully flat diffuse upper and different exact corner replacement
already proved above finish the theorem.

\paragraph{The lower bound of Remark~\ref{rem:general-houghton}.} By Theorem~\ref{thm:extraction} with $\lambda=2$ and Proposition~\ref{prop:purity}, some round of a rank-rate
service with $\epsilon\le\frac1{16}$ has distance at most $\frac1{16}$ and $\ell+a\le\frac{32}{15}(3Q+7+\log n)/n$. The one-round theorem~\cite[Theorem~5.5]{DouglasMghirbiB} gives
$Q+\frac18\ge S(\rho)+\frac18\ge c_*(\ell+a)^{-\beta_H}$. If $Q\ge7+\log n$, then $\ell+a\le\frac{128}{15}Q/n$ and $Q^{1+\beta_H}\ge c\,n^{\beta_H}$. If $Q<7+\log n$, then
$Q+\frac18\ge c_*(15n/(128(7+\log n)))^{\beta_H}$, which is impossible for large $n$. So $\Qs_{\Phi_H}(n,\epsilon)\ge c\,n^{\beta_H/(1+\beta_H)}$ for $n\ge n_0$.
